\documentclass[10pt, a4paper]{article}

% - Packages ---
\usepackage[top=2cm, bottom=2cm, left=2cm, right=2cm]{geometry}
\usepackage{graphicx} % For images
\usepackage{amsmath}  % For math/equations
\usepackage{hyperref} % For clickable links
\usepackage{caption}  % For styling captions
\usepackage{titlesec} % For controlling section title sizes
\usepackage{float}
\usepackage{xcolor} % Required for defining and using colors
\usepackage{hyperref} % Keep this as your last \usepackage
\usepackage{color}
\usepackage{soul}

% ברירת המחדל של המרקור היא צהוב. אם תרצי לשנות למשל לטורקיז/תכלת, הוסיפי גם את השורה הזו:
% \sethlcolor{cyan}

\hypersetup{
    colorlinks=true, % Removes the default boxes and colors the text instead
    linkcolor=blue,  % Color for internal links (e.g., Figure and Table references)
    citecolor=blue,  % Color for bibliographical citations
    urlcolor=blue    % Color for external web URLs
}
% - Document Formatting ---
\setlength{\parindent}{0pt}
\setlength{\parskip}{0.5em}
\captionsetup{font=scriptsize, labelfont=bf}

% - Customize Title Sizes ---
\titleformat{\section}
  {\large\bfseries}{\thesection}{1em}{}
\titleformat{\subsection}
  {\normalsize\bfseries}{\thesubsection}{1em}{}

\usepackage{fancyhdr}
\pagestyle{fancy}
\fancyhf{} % מנקה את הגדרות ברירת המחדל
\fancyhead[L]{\footnotesize Weekly Summary \#6 \today} % יופיע למעלה בשמאל
\fancyhead[R]{\footnotesize Rachel Broner $\vert$ Miri Adler's Lab} % יופיע למעלה בימין
\fancyfoot[C]{\thepage} % מספר עמוד למטה באמצע

% \usepackage{mathptmx}

\usepackage{setspace}
\onehalfspacing % For 1.5 line spacing


\begin{document}

\begin{center}
    {\Large \textbf{Weekly Summary \#6 $\vert$ Association Rule Mining}} \\[0.5em]
    {\large Part 3 $\vert$ Which cell type sits next to which} \\[0.5em]
    \textit{GVHD Analysis} $\vert$ Week of: \today
\end{center}
\vspace{0.2cm}
\hrule
\vspace{0.5cm}

\textbf{Muscle issue:} How much muscularis a biopsy caught
differs by stage, and in opposite directions in the two organs: FOVs containing muscle rise
with stage in the duodenum (58\%, 66\%, 68\%) and fall in the colon (80\%, 33\%, 52\%). Any
result involving Muscle therefore mixes biology with what the biopsy happened to include,
so every figure below that touches it is shown twice - all FOVs, then only those that
sampled the muscularis. Restricting this way loses no patient from either severe group.

% ---------------------------------------------------------
\section{Claim 1 - the direction of a rule is anatomical}

\textit{A $\rightarrow$ B} and \textit{B $\rightarrow$ A} are different measurements. \textbf{59\%} of mined rules exist in one direction only.

\textbf{In one image, two squares are filled and their mirrors are empty.} Paneth cells
have epithelium beside them; smooth muscle has endothelium beside it. Neither reverse rule
was mined - which is what the anatomy predicts, since Paneth cells are under 1\% of the
tissue and almost no epithelial cell has one nearby.

\begin{figure}[H]
    \centering
    \includegraphics[width=0.78\textwidth]{summary_downloads/f1_one_fov_control.pdf}
    \caption{\textbf{Reading the grid:} Center cell down the side, neighbor cell across the
    bottom, one dot per mined rule. Color is lift on a log2 spacing, so 0.5 and 2 sit
    equally far from white. Dot size is the share of FOVs carrying the rule; the ring is
    that share among FOVs holding both cell types, so a gap means the cells are often
    absent. Here: one healthy duodenal FOV (Control\_01\_FOV\_1), axes trimmed to the cell
    types it uses.}
    \label{fig:fov1}
\end{figure}

\textbf{Across the cohort the asymmetry holds:}
Paneth $\rightarrow$ Epithelial fires in 12 of 29 eligible control FOVs; its mirror in
\textbf{0 of 29}. SMV $\rightarrow$ Endothelial sits well above chance at every stage while
its mirror reaches 1 of 30 controls.

\begin{figure}[H]
    \centering
    \includegraphics[width=\textwidth]{summary_downloads/f2_direction_pairs.pdf}
    \caption{\textbf{Reading a dot panel:} One cell pair per panel; each dot is one FOV
    carrying the rule, the bar is the median, and the fraction under each stage is FOVs with
    the rule out of FOVs holding both cell types. The two right panels are the mirrors of
    the two left ones.}
    \label{fig:dir}
\end{figure}

\textbf{Pooled, the grid is visibly not symmetric about its diagonal,} with the
asymmetry the right way round.

\begin{figure}[H]
    \centering
    \includegraphics[width=0.78\textwidth]{summary_downloads/f3_group_control.pdf}
    \caption{Healthy duodenum, 31 FOVs pooled, axes trimmed as before.}
    \label{fig:group1}
\end{figure}

% ---------------------------------------------------------
\section{Claim 2 - the rules reconstruct the layers of the gut}

The gut is built in layers - epithelium, lamina propria, muscularis - and nothing about
this was given to the mining.

\textbf{Folded into lineages, grouped cells shows the layering.} The diagonal is warm
(a layer sits with itself) and the epithelial row is cold against everything outside the
epithelium. Epithelium against muscle is the strongest avoidance anywhere in the data.

\begin{figure}[H]
    \centering
    \includegraphics[width=0.7\textwidth]{summary_downloads/f4a_layers_one.pdf}
    \caption{Healthy duodenum, the 31 cell types folded into their lineages; Muscle and
    Brunner glands keep their own row and column, forming a lineage with nothing else.}
    \label{fig:layers1}
\end{figure}

\textbf{The same picture repeats in both organs and at every stage.} Epithelium keeps apart
from immune cells (lift 0.36--0.62) and from stroma (0.26--0.48); stroma sits with stroma
(1.33--1.54); epithelium against muscle stays at 0.03--0.09. This is the stereotypic
organisation Azulay \textit{et al.} describe for healthy gut, recovered from rules alone,
and it survives aGVHD intact.

\begin{figure}[H]
    \centering
    \includegraphics[width=\textwidth]{summary_downloads/f4b_layers_all.pdf}
    \caption{All six organ $\times$ stage groups, folded as in Fig.~\ref{fig:layers1}.}
    \label{fig:layers}
\end{figure}

\textbf{Dropping the FOVs with no muscularis changes nothing but the Muscle row.} The
diagonal, the epithelial row and the stromal block are the same picture, so the layering is
not an artefact of which biopsies caught muscle.

\begin{figure}[H]
    \centering
    \includegraphics[width=\textwidth]{summary_downloads/f4c_layers_muscularis.pdf}
    \caption{Fig.~\ref{fig:layers} restricted to FOVs that sampled the muscularis.}
    \label{fig:layersmusc}
\end{figure}

\textbf{Unfolded, each pair sits at the same height in every stage - the boundary does not
loosen.} Tested one rule at a time: of 24 across-layer rules in the duodenum, none moves
toward chance at BH FDR $<0.10$, and of 38 in the colon only 2 do. Several tighten instead
--- Goblet $\rightarrow$ Plasma goes $-0.82 \rightarrow -1.76$ in the duodenum and
$-0.90 \rightarrow -1.72$ in the colon. The paper's ``less organised'' is about villus and crypt
\emph{shape}; a rule only asks which cell types end up within $25\,\mu$m of each other, and
a crypt can lose its shape while its cells keep their neighbours.

\begin{figure}[H]
    \centering
    \includegraphics[width=0.90\textwidth]{summary_downloads/f5_layer_pairs.pdf}
    \caption{Colon. Three pairs across the layer boundary, drawn as in
    Fig.~\ref{fig:dir}.}
    \label{fig:layerpairs}
\end{figure}

\textbf{Epithelial $\rightarrow$ Muscle is the one pair the sampling could have invented,
and it survives.} On muscularis FOVs only it stays far below chance at every stage; the
other two pairs are unaffected, as they involve no muscle.

\begin{figure}[H]
    \centering
    \includegraphics[width=0.90\textwidth]{summary_downloads/f5b_layer_pairs_muscularis.pdf}
    \caption{Fig.~\ref{fig:layerpairs} restricted to FOVs that sampled the muscularis.}
    \label{fig:layerpairsmusc}
\end{figure}

% ---------------------------------------------------------
\section{Claim 3 - the paper's findings appear as rules arriving and leaving}

\textbf{One diseased image already shows two of the paper's epithelial findings.} Against
Fig.~\ref{fig:fov1}, the Paneth row has gone and an Endocrine row has appeared.

\begin{figure}[H]
    \centering
    \includegraphics[width=0.78\textwidth]{summary_downloads/f6_one_fov_severe.pdf}
    \caption{One severe duodenal FOV (GVHD\_14\_FOV\_4), drawn as Fig.~\ref{fig:fov1}.}
    \label{fig:fov2}
\end{figure}

\textbf{Across the cohort, Paneth rules vanish and Endocrine rules appear - but the dots
do not move up or down, only their number changes.} Paneth $\rightarrow$ Epithelial holds
through mild disease (12/29 then 37/91 of eligible FOVs) and drops in severe (7/30).
Endocrine $\rightarrow$ Epithelial is absent from healthy duodenum entirely (0/21) and
present only with disease (22/87 mild, 7/35 severe).

\begin{figure}[H]
    \centering
    \includegraphics[width=0.72\textwidth]{summary_downloads/f7_paper_duodenum.pdf}
    \caption{Duodenum, drawn as in Fig.~\ref{fig:dir}. The vertical position of the dots is
    the rule's strength; the count beneath is how many FOVs carry it.}
    \label{fig:paper1}
\end{figure}

\textbf{The cells themselves moved, which is why the rules did.} Paneth cells fall from
0.9\% of an FOV to 0.6\% and from 29 of 31 FOVs to 30 of 41; endocrine cells rise from
0.1\% to 0.5\% and from 21 of 31 to 35 of 41. So these rules track the paper's
\emph{compositional} findings, not a rearrangement.

\begin{figure}[H]
    \centering
    \includegraphics[width=0.72\textwidth]{summary_downloads/f7b_cells_duodenum.pdf}
    \caption{The same FOVs, now showing how much of each one the cell type makes up. Same
    layout as Fig.~\ref{fig:paper1}, so the two can be read against each other.}
    \label{fig:cells1}
\end{figure}

\textbf{Both rules are sharper against the pathological score than the clinical one.} Grouped
by pathology, Paneth $\rightarrow$ Epithelial falls further in severe disease (2/24, 8\%,
against 7/30, 23\% clinically) and Endocrine $\rightarrow$ Epithelial keeps rising
(0\%, 22\%, 30\%) instead of levelling off (0\%, 25\%, 20\%). Both are histological changes
and the Lerner pathological score is the histological readout, while the Glucksberg clinical
score grades symptoms - the paper notes the two disagree in six of its patients. That the
epithelial rules track the tissue score more closely is a small point in the method's favour:
they are responding to damage rather than to how ill the patient felt.

\begin{figure}[H]
    \centering
    \includegraphics[width=0.72\textwidth]{summary_downloads/f7c_paper_duodenum_pathological.pdf}
    \caption{Fig.~\ref{fig:paper1} regrouped by the pathological score. Everything else in
    this note uses the clinical score, as the paper does.}
    \label{fig:paper1path}
\end{figure}

\textbf{The colon repeats the pattern in both directions.} Neutrophil $\rightarrow$
Epithelial is absent from control and mild disease and appears only in severe (4/52);
CD4T $\rightarrow$ Plasma runs the other way, from 14/25 down to 2/50.

\begin{figure}[H]
    \centering
    \includegraphics[width=0.72\textwidth]{summary_downloads/f8_paper_colon.pdf}
    \caption{Colon, drawn as in Fig.~\ref{fig:dir}.}
    \label{fig:paper2}
\end{figure}

\textbf{Again the cells explain it}: neutrophils spread with severity while plasma and CD4 T
cells fall away - the neutrophil enrichment and plasma/CD4T depletion the paper reports.

\begin{figure}[H]
    \centering
    \includegraphics[width=\textwidth]{summary_downloads/f8b_cells_colon.pdf}
    \caption{Colon cell abundance, drawn as in Fig.~\ref{fig:cells1}.}
    \label{fig:cells2}
\end{figure}

Requiring at least 20 cells of each type does not change this: Paneth $\rightarrow$
Epithelial still falls from 77\% to 54\% of FOVs, but no pair's \emph{lift} differs between
control and severe after correction. These figures confirm the paper's findings in the rule
table; they do not rediscover them independently.

\textbf{All six groups at once.} The Goblet rows say it most clearly: small dots inside
large rings, so the rules hold wherever goblet cells remain and there are simply fewer such
FOVs.

\begin{figure}[H]
    \centering
    \includegraphics[width=\textwidth]{summary_downloads/f9_all_groups.pdf}
    \caption{All six groups on a shared axis, restricted to the 14 cell types carrying a
    rule in over 25\% of the FOVs of at least three groups.}
    \label{fig:all}
\end{figure}

\textbf{On muscularis FOVs only, the epithelial changes are unchanged and the Muscle rows
move.} The Paneth, Endocrine and Goblet rows read the same, so the findings above do not
depend on which biopsies caught muscle.

\begin{figure}[H]
    \centering
    \includegraphics[width=\textwidth]{summary_downloads/f10_muscularis.pdf}
    \caption{Fig.~\ref{fig:all} restricted to FOVs that sampled the muscularis.}
    \label{fig:muscularis}
\end{figure}

% ---------------------------------------------------------
\section{What the method cannot answer}

\textbf{The grid is sparse.} A median FOV yields 30 rules out of 930 possible pairs, so an
untrimmed grid is about 97\% empty. Trimming to the cell types actually used lifts a pooled
group grid from 5\% to 27\% filled; the shared axis is kept only where grids must be compared.

\textbf{Abundant cell types are handicapped.} A cell type in nearly every neighborhood has
a lift near 1 by construction and cannot clear the threshold. This puts the paper's
intraepithelial-lymphocyte result out of reach: the signs are right at every stage
(CD8T $\rightarrow$ Epithelial above 1, CD4T $\rightarrow$ Epithelial below 1) but the
rules clear threshold in only 1--3 FOVs.

\textbf{One FOV is coarser than one crypt.} The paper's local-microenvironment result
(their Fig.~6) is measured in a ring around each crypt, where an infiltrated crypt can sit
beside an untouched one. A rule is mined over a whole FOV, which averages the two together.
Within-FOV patchiness would need mining at crypt level - the natural next step.

\end{document}
